% GENERATED FILE. Edit canonical lessons, then run npm run generate:books.
% canonical-source-hash: fnv1a64:d3dbe97f6345db6d
% canonical-lessons: TA-C256-errukkol, TA-C256-pinparru, TA-C256-tavaravitu, TA-C256-ali, TA-C256-accitu

\chapter{To accept, To follow, To miss, To erase, To print}
\label{ch:ta-to-accept-to-follow-to-miss-to-erase-to-print}
\hlchaptermodality{256}

\begin{chapteropening}
\textbf{By the end of this chapter:} \emph{I can say to accept, to follow, to miss, to erase and to print.}

Complete the last lesson of chapter 256: I can say to accept, to follow, to miss, to erase and to print.
\end{chapteropening}

\section[\texorpdfstring{\ta{ஏற்றுக்கொள்} (ēṟṟukkoḷ)}{ஏற்றுக்கொள் (ēṟṟukkoḷ)}]{\ta{ஏற்றுக்கொள்} (ēṟṟukkoḷ) — to accept}

\label{lesson:TA-C256-errukkol}



\begin{quote}
Before the new one: say the Tamil for to wake, then the Tamil for to brush one's teeth.
\end{quote}

\subsection*{You'll want to know: \ta{ஏற்றுக்கொள்}}
\textbf{\ta{ஏற்றுக்கொள்}} — \emph{ēṟṟukkoḷ} — \textquotedblleft{}to accept\textquotedblright{}.

\begin{grammarlens}[title={What you've built}]
One more word for this chapter.
\end{grammarlens}

\subsection*{Your turn}
\begin{itemize}
\raggedright
  \item \emph{Say it:} \emph{ēṟṟukkoḷ}
  \item \emph{Say it:} \emph{ēṟṟukkoḷ}, once more
  \item \emph{Say it:} say \emph{ēṟṟukkoḷ}
  \item \emph{Recall:} say the Tamil for to wake, then the Tamil for to brush one's teeth, then say \emph{ēṟṟukkoḷ} again
\end{itemize}

\begin{tcolorbox}[breakable,colback=teal!4,colframe=teal!35!black,title={Before you move on}]
What does \ta{ஏற்றுக்கொள்} mean? (\textquotedblleft{}To accept\textquotedblright{}.) Say it once more.
\end{tcolorbox}
\section[\texorpdfstring{\ta{பின்பற்று} (piṉpaṟṟu)}{பின்பற்று (piṉpaṟṟu)}]{\ta{பின்பற்று} (piṉpaṟṟu) — to follow (rules, directions)}

\label{lesson:TA-C256-pinparru}



\begin{quote}
Before the new one: say the Tamil for to brush one's teeth, then the Tamil for to accept.
\end{quote}

\subsection*{You'll want to know: \ta{பின்பற்று}}
\textbf{\ta{பின்பற்று}} — \emph{piṉpaṟṟu} — \textquotedblleft{}to follow (rules, directions)\textquotedblright{}.

\begin{grammarlens}[title={What you've built}]
One more word for this chapter.
\end{grammarlens}

\subsection*{Your turn}
\begin{itemize}
\raggedright
  \item \emph{Say it:} \emph{piṉpaṟṟu}
  \item \emph{Say it:} \emph{piṉpaṟṟu}, once more
  \item \emph{Say it:} say \emph{piṉpaṟṟu}
  \item \emph{Recall:} say the Tamil for to brush one's teeth, then the Tamil for to accept, then say \emph{piṉpaṟṟu} again
\end{itemize}

\begin{tcolorbox}[breakable,colback=teal!4,colframe=teal!35!black,title={Before you move on}]
What does \ta{பின்பற்று} mean? (\textquotedblleft{}To follow (rules, directions)\textquotedblright{}.) Say it once more.
\end{tcolorbox}
\section[\texorpdfstring{\ta{தவறவிடு} (tavaṟaviṭu)}{தவறவிடு (tavaṟaviṭu)}]{\ta{தவறவிடு} (tavaṟaviṭu) — to miss (a bus, a chance)}

\label{lesson:TA-C256-tavaravitu}



\begin{quote}
Before the new one: say the Tamil for to accept, then the Tamil for to follow.
\end{quote}

\subsection*{You'll want to know: \ta{தவறவிடு}}
\textbf{\ta{தவறவிடு}} — \emph{tavaṟaviṭu} — \textquotedblleft{}to miss (a bus, a chance)\textquotedblright{}.

\begin{grammarlens}[title={What you've built}]
One more word for this chapter.
\end{grammarlens}

\subsection*{Your turn}
\begin{itemize}
\raggedright
  \item \emph{Say it:} \emph{tavaṟaviṭu}
  \item \emph{Say it:} \emph{tavaṟaviṭu}, once more
  \item \emph{Say it:} say \emph{tavaṟaviṭu}
  \item \emph{Recall:} say the Tamil for to accept, then the Tamil for to follow, then say \emph{tavaṟaviṭu} again
\end{itemize}

\begin{tcolorbox}[breakable,colback=teal!4,colframe=teal!35!black,title={Before you move on}]
What does \ta{தவறவிடு} mean? (\textquotedblleft{}To miss (a bus, a chance)\textquotedblright{}.) Say it once more.
\end{tcolorbox}
\section[\texorpdfstring{\ta{அழி} (aḻi)}{அழி (aḻi)}]{\ta{அழி} (aḻi) — to erase, to delete}

\label{lesson:TA-C256-ali}



\begin{quote}
Before the new one: say the Tamil for to follow, then the Tamil for to miss.
\end{quote}

\subsection*{You'll want to know: \ta{அழி}}
\textbf{\ta{அழி}} — \emph{aḻi} — \textquotedblleft{}to erase, to delete\textquotedblright{}.

\begin{grammarlens}[title={What you've built}]
One more word for this chapter.
\end{grammarlens}

\subsection*{Your turn}
\begin{itemize}
\raggedright
  \item \emph{Say it:} \emph{aḻi}
  \item \emph{Say it:} \emph{aḻi}, once more
  \item \emph{Say it:} say \emph{aḻi}
  \item \emph{Recall:} say the Tamil for to follow, then the Tamil for to miss, then say \emph{aḻi} again
\end{itemize}

\begin{tcolorbox}[breakable,colback=teal!4,colframe=teal!35!black,title={Before you move on}]
What does \ta{அழி} mean? (\textquotedblleft{}To erase, to delete\textquotedblright{}.) Say it once more.
\end{tcolorbox}
\section[\texorpdfstring{\ta{அச்சிடு} (acciṭu)}{அச்சிடு (acciṭu)}]{\ta{அச்சிடு} (acciṭu) — to print}

\label{lesson:TA-C256-accitu}



\begin{quote}
Before the new one: say the Tamil for to miss, then the Tamil for to erase.
\end{quote}

\subsection*{You'll want to know: \ta{அச்சிடு}}
\textbf{\ta{அச்சிடு}} — \emph{acciṭu} — \textquotedblleft{}to print\textquotedblright{}.

\begin{grammarlens}[title={What you've built}]
One more word for this chapter.
\end{grammarlens}

\subsection*{Your turn}
\begin{itemize}
\raggedright
  \item \emph{Say it:} \emph{acciṭu}
  \item \emph{Say it:} \emph{acciṭu}, once more
  \item \emph{Say it:} say \emph{acciṭu}
  \item \emph{Recall:} say the Tamil for to miss, then the Tamil for to erase, then say \emph{acciṭu} again
\end{itemize}

\begin{tcolorbox}[breakable,colback=teal!4,colframe=teal!35!black,title={Before you move on}]
What does \ta{அச்சிடு} mean? (\textquotedblleft{}To print\textquotedblright{}.) Say it once more.
\end{tcolorbox}
